\documentclass[12pt]{article}
\usepackage[utf8]{inputenc}
\usepackage[T1]{fontenc}
\usepackage{hyperref}
\title{Raetic Language Resources for K-12 Curriculum Development}
\author{Compiled by Claude Sonnet 5.5 (claude-sonnet-5-5) for Andr\'as Kornai}
\date{}
\begin{document}
\maketitle

\section{Introduction}
Raetic (ISO 639-3: xrr; Glottocode raet1238) is an ancient, non-Indo-European language of the eastern Alpine region, known only from inscriptions dated roughly to the 6th to 1st centuries BC. Findspots are mainly in Trentino, South and North Tyrol, with some in the Veneto, Graub\"unden and Slovenia. It is generally related to the Tyrsenian family (with Etruscan and Lemnian), though the glottocode record here lists it as unclassifiable. There are no speakers, and no school curriculum in Raetic was found. Any curriculum would be a heritage or history unit, not instruction in the language.

\section{Primary Educational Resources}
None found. No primers, textbooks or school programmes for Raetic appeared in the search results.

\section{Linguistic and Reference Materials}
\begin{itemize}
\item Thesaurus Inscriptionum Raeticarum (TIR), University of Vienna, the main digital corpus edition: \url{https://tir.univie.ac.at/}
\item TIR wiki page on Raetica: \url{https://tir.univie.ac.at/wiki/Raetica}
\item TIR inscription category listing (about 389 inscriptions, 112 classed as certainly Raetic): \url{https://tir.univie.ac.at/wiki/Category:Inscription}
\item FWF project page for the TIR project P 25495: \url{https://www.fwf.ac.at/en/research-radar/10.55776/P25495}
\item Digital Classicist wiki entry on the TIR: \url{https://wiki.digitalclassicist.org/Thesaurus_Inscriptionum_Raeticarum}
\item Overview page on Raetic (TLLM wiki): \url{https://episteme.tllm.fr/wiki/Raetic}
\item PRIN Italia Antica project page: \url{https://www.prin-italia-antica.unifi.it/p155.html}
\end{itemize}
The results also cite Salomon (2018), ``Thesaurus Inscriptionum Raeticarum -- New readings and inscriptions'', \emph{Die Sprache}; no URL for it was seen.

\section{Religious and Community Resources}
None found.

\section{Digital and Audio Resources}
The TIR website is the principal digital resource. Its semantic-wiki software is described at \url{https://semantic-mediawiki.org/wiki/Thesaurus_Inscriptionum_Raeticarum_(TIR)}. No audio exists.

\section{Curriculum Development Context}
Because Raetic is attested only in short inscriptions, a K-12 treatment would be a cultural and historical unit on the pre-Roman Alps, using the TIR inscriptions as primary sources. Alphabets, inscription reading and archaeology would suit upper-secondary classes. Local schools in the Tyrol and Trentino regions would be natural partners; no specific programme was found.

\end{document}